% !TEX program = xelatex
% Paper 2 in Molecular Informatics (Wiley New Journal Design v5) format.
% COMPILE ON OVERLEAF using the "Wiley NJD-v5" template (XeLaTeX), which
% supplies WileyNJDv5.cls and the wileyNJD-*.bst files. This file replaces
% that template's main .tex. Figures (fig2/3/5b/8/9, fig_graphical_abstract)
% must be uploaded alongside.
\documentclass[STIX1COL]{WileyNJDv5}

% --- math / tables / lists used by the body ---
\usepackage{amsmath,amssymb,amsthm}
\usepackage{mathtools}
\usepackage{booktabs}
\usepackage{array}
\usepackage{enumitem}
\newtheorem{theorem}{Theorem}
\theoremstyle{definition}
\newtheorem{remark}[theorem]{Remark}
\newcommand{\bid}{\mathcal{F}_{\mathrm{BID}}}

\articletype{Research Article}

\title{Orthogonality Screening of Topological Indices for QSPR Modeling: A Structural and Empirical Redundancy Analysis}

\author[1]{Ganesh Shiwakoti}
\author[2]{Advik Natarajan}
\author[3]{Parik Chalise}
\author[2]{Michael Arockiaraj}

\authormark{Shiwakoti \textsc{et al.}}
\titlemark{Orthogonality Screening of Topological Indices}

\address[1]{Department of Mathematics and Statistics, Florida Atlantic University, Florida, USA}
\address[2]{Department of Mathematics, Loyola College, Chennai, India}
\address[3]{Department of Applied Mathematics and Statistics, Johns Hopkins University, Baltimore, USA}

\corres{Ganesh Shiwakoti, Department of Mathematics and Statistics, Florida Atlantic University, USA. \email{ganshiwakoti@gmail.com}}

\abstract[Abstract]{The chemical-graph-theory literature contains hundreds of scalar
topological indices, and the usual test for whether a newly
proposed index carries independent information --- a pairwise
correlation below $0.95$ against a classical baseline --- is
empirically too permissive. We present an open-source screening
pipeline that combines pairwise correlation, principal-component
effective rank, variance inflation, and partial correlation with
the modelling target into a single pre-flight check, complemented
by a structural result (the \emph{Edge-Degree-Pair Basis}) showing
that every bond-incident-degree index on hydrogen-suppressed
molecular graphs of maximum degree four lies in a vector space of
dimension at most ten. Applied to a three-parameter parametric
index family and to twenty alternative candidate indices across
four MoleculeNet benchmarks, the screen shows that most candidates
clear the pairwise test yet add little independent signal once
partial correlation is also required; a cross-validated
machine-learning benchmark confirms that the pruned index set
preserves predictive performance at several-fold fewer features.
We conclude that orthogonality is necessary but not sufficient for
a new index to be useful, and provide the pipeline as a reusable
standard --- noting that the diagnostics are linear and that the
screen is a pre-flight check, not a hard feature-selection filter.}

\keywords{Topological index, orthogonality screening, redundancy analysis, QSPR, principal component analysis, partial correlation, bond-incident-degree (BID) family}

\begin{document}

\maketitle

% --- Graphical abstract / TOC entry (Molecular Informatics) ---
% \graphicalabstract{fig_graphical_abstract.pdf}{Only 0--4 of 20 candidate
% topological indices carry independent QSPR signal once partial correlation
% is required; the bond-incident-degree family lies in a 10-dimensional space.}

\section{Introduction}
% =====================================================================

The vocabulary of scalar topological indices used in QSPR/QSAR
modeling has grown over six decades from a handful of
proposals \cite{Wiener1947, Randic1975, Gutman1972} to several
hundred, and the resulting proliferation has been criticised
within mathematical chemistry itself~\cite{Gutman2013Degree}.
Many indices are interrelated by exact identity (e.g.\
the harmonic index $H$ is $2$ times the special case
$\beta=-1$ of the general sum-connectivity index $\chi_\beta$), and
in practice the empirical Pearson correlations between proposed
indices on a fixed dataset can exceed $|r|=0.95$ for the majority
of pairs.

This raises a recurring methodological question. \emph{What does a
newly proposed scalar topological index need to demonstrate to be
treated as a meaningful contribution to QSPR modeling?} Three
weak forms of justification appear frequently in the literature:
\begin{enumerate}[label=(\roman*)]
\item Strong correlation with a chemical property on a small dataset
(e.g.\ $18$ octanes, $75$ decanes), often without reference to a
classical-baseline comparison.
\item A closed-form value on standard graph classes, framed as
mathematical novelty.
\item Parametric generality (a new parameter that subsumes the
classical case at one specific parameter value).
\end{enumerate}
None of these by themselves is sufficient to demonstrate
\emph{independent} predictive value beyond the classical baseline.

This paper proposes an explicit pre-flight check for the
``does it carry independent information'' question, framed
operationally so that any author can apply it.

\subsection{Contributions}

\begin{enumerate}[label=(\arabic*)]
\item An open-source orthogonality-screening pipeline
(Section~\ref{sec:protocol}) covering correlation matrix,
PCA, VIF, and partial correlation with the QSPR target, using
a fixed redundancy threshold $|r|\ge 0.95$.
\item Replication across four MoleculeNet benchmarks: three
regression tasks (ESOL, FreeSolv, Lipophilicity) and one binary
classification task (BBBP)
(Section~\ref{sec:datasets}).
\item A structural observation (Section~\ref{sec:bid}, the
\emph{Edge-Degree-Pair Basis}) giving a dimension-$10$ ceiling on
the BID family on $\mathcal{G}_4$ (hydrogen-suppressed molecular
graphs).
\item Two case studies: the parametric Wuzi index family
(Section~\ref{sec:wuzi_case}, $300$ grid points) and twenty
alternative-family candidate indices
(Section~\ref{sec:novel_case}).
\item Explicit positioning of orthogonality as \emph{necessary but
not sufficient} for QSPR usefulness, with empirical evidence
(Section~\ref{sec:discussion}).
\end{enumerate}

% =====================================================================
\section{Related work}\label{sec:related}
% =====================================================================

\subsection{Relation to the companion graph-theory paper}\label{sec:relation_companion}

The mathematical-chemistry properties of the Wuzi index family ---
its definition, closed-form values on standard graph classes,
bracket and ratio bounds, Cauchy--Schwarz and Jensen-type bounds
in $M_1$, $M_2$, $R$, $\chi$, and $SO$, the empirical
extremal-graph pattern (including the spider / double-spider
structural conjecture at the mixed-sign triple
$(\alpha, \beta, \gamma) = (-1, -1, 1)$ verified through tree
order $n = 20$), and sensitivity-analysis benchmarks on $18$
octanes, $106$ trees of order $10$, and $75$ decane isomers ---
are developed in a companion graph-theory
paper~\cite{Paper1_Wuzi}. The present paper is focused on the
screening pipeline as a reusable methodology and uses the Wuzi
family only as one screened candidate (Section~\ref{sec:wuzi_case}),
re-deriving no Wuzi-specific mathematical results. A reader
interested only in the screening methodology can treat the Wuzi
family as a black-box parametric BID family throughout.

\subsection{Orthogonality and redundancy of topological indices}

The recognition that many topological indices are mutually highly
correlated is not new. \cite{MovahediGutmanRedzepovicFurtula2026}
Section~5.2 includes an intercorrelation analysis among
Sombor-family variants showing pairwise $r$-values above $0.9$ on
the octane data. Other intercorrelation observations appear
sporadically in the BID-bounds literature. The present work
extends this style of observation in three ways: (i)~by fixing a
threshold $|r|\ge 0.95$ and applying it uniformly across datasets;
(ii)~by replicating across three regression QSPR benchmarks and
one binary classification benchmark rather than a single small
dataset; and (iii)~by combining
PCA, partial correlation, and VIF with the pairwise correlation
analysis to assess different aspects of redundancy.

\subsection{The bond-incident-degree (BID) family}

\emph{Bond-incident-degree} (BID) indices are functionals of the
form $I_f(G)=\sum_{uv\in E(G)} f(d_u,d_v)$ for a symmetric
function $f$. This unifies the bulk of the classical degree-based
topological-index literature: Zagreb $M_1,M_2$, modified Zagreb,
forgotten $F$, Randi\'c $R$ and general $R_\alpha$,
sum-connectivity $\chi$ and general $\chi_\beta$, harmonic $H$,
the geometric--arithmetic index $GA$ of Vuki\v{c}evi\'c and
Furtula~\cite{VukicevicFurtula2009GA}, arithmetic-geometric $AG$,
atom-bond connectivity $ABC$, augmented Zagreb $AZI$, Sombor $SO$
and its variants, the Albertson irregularity index
$\mathrm{Alb}$~\cite{Albertson1997}, sigma $\sigma$, and reduced
Zagreb $\mathrm{redM}_1$. The observation that BID
indices are linear functionals of the edge-degree-pair counts
$m_{ij}$ is standard~\cite{BorovicaninDasFurtulaGutman2017Zagreb}; what we
contribute is the explicit dimension bound on $\mathcal{G}_4$ and
the redundancy-screening consequence.

\subsection{QSPR benchmarks and reproducibility}

The QSPR/QSAR benchmarks used here are drawn from the MoleculeNet
collection~\cite{MoleculeNet}, a standard reference for molecular
property prediction: three regression benchmarks (ESOL, FreeSolv,
Lipophilicity) and one binary classification benchmark (BBBP). We
use the versions distributed via MoleculeNet without resplitting;
SMILES are processed through RDKit~\cite{RDKit} and converted to
hydrogen-suppressed molecular graphs through a thin NetworkX
wrapper.

\subsection{Relation to generic feature-selection methods}\label{sec:feature_selection}

The broader cheminformatics descriptor-redundancy literature is
documented in the Todeschini--Consonni
handbook~\cite{TodeschiniConsonniHandbook}, which catalogues the
classical molecular-descriptor families and the standard pairwise
and multivariate redundancy diagnostics. Outside cheminformatics,
generic feature-selection methods --- the minimum-redundancy
maximum-relevance criterion (mRMR)~\cite{PengLongDing2005},
the Boruta wrapper~\cite{KursaRudnicki2010}, and embedded
methods such as LASSO~\cite{Tibshirani1996} --- select features
from an arbitrary descriptor matrix conditional on a fixed target
and model class.

The present pipeline is positioned at a different layer in the
modelling stack. The pairwise / PCA / VIF diagnostics are
\emph{target-agnostic}; the partial-correlation step is
\emph{target-aware but model-agnostic}; and the combined diagnostic
is designed as a \emph{pre-flight check for a single proposed new
topological index} against a fixed classical baseline, not as a
competitor for end-to-end QSPR feature selection. Conceptually,
the two layers compose rather than compete: the screening pipeline
gates which indices enter the descriptor matrix, after which
mRMR / LASSO / Boruta / similar methods could be applied
downstream to the screened set against the modelling task at hand.

The present paper includes both a mixed-model sanity check
against LASSO, ElasticNet, mRMR, Boruta, and $\ell_1$-logistic
feature selection in Section~\ref{sub:lasso_comparator} and a
matched-model-class pipeline-then-LASSO / $\ell_1$-logistic
benchmark in Section~\ref{sub:pipeline_then_lasso}. Broader
matched comparisons involving nonlinear or tree-ensemble
downstream models (e.g.\ pipeline-then-XGBoost,
pipeline-then-RandomForest with Boruta- or mRMR-style selection)
remain future work.

% =====================================================================
\section{Datasets and implementation}\label{sec:datasets}
% =====================================================================

\begin{table}[h]
\centering
\caption{QSPR/QSAR benchmark datasets used in this study. The $n$ column reports the number of molecules that survive RDKit parsing (largest connected fragment, $\ge 2$ heavy atoms).}
\label{tab:datasets}
\begin{tabular}{lllrl}
\toprule
Name & Target property & Task & $n$ & Source\\
\midrule
ESOL          & $\log S$ (aqueous solubility, mol/L)         & regression     & $1\,127$ & \cite{Delaney2004ESOL} via \cite{MoleculeNet}\\
FreeSolv      & $\Delta G_{\mathrm{hyd}}$ (hydration free energy, kcal/mol) & regression & $639$    & \cite{Mobley2014FreeSolv} via \cite{MoleculeNet}\\
Lipophilicity & $\log D$ at pH 7.4                            & regression     & $4\,200$ & \cite{LipophilicityChEMBL} via \cite{MoleculeNet}\\
BBBP          & blood-brain barrier penetration               & classification & $2\,039$ & \cite{Martins2012BBBP} via \cite{MoleculeNet}\\
\bottomrule
\end{tabular}
\end{table}

Table~\ref{tab:datasets} summarises the four benchmarks; the
$30$-index classical baseline (Table~\ref{tab:baseline}) is
implemented in the reference codebase
(\texttt{src/standard\_indices.py}) on top of
NetworkX, RDKit~\cite{RDKit}, and the scikit-learn machine-learning
toolkit~\cite{Pedregosa2011sklearn}. Each index admits a
$\le 20$-line Python implementation, and the full $30$-index
computation for a single graph takes a few milliseconds. The
reference implementation is open-source under MIT license.

\noindent\emph{On the baseline composition.} The $30$-index set
follows the practitioner-standard cheminformatics
baseline~\cite{TodeschiniConsonniHandbook}: degree-based,
distance-based, spectral, and other-family indices in the
proportions in which they are actually used in published QSPR
modelling. Eighteen of the thirty are degree-based BID indices;
this is deliberate. Theorem~\ref{thm:bid} shows that the BID
subspace on $\mathcal{G}_4$ has dimension at most $10$, so the
baseline needs enough BID indices to span (or come close to
spanning) that subspace in order for residual partial
correlations of BID candidates to be informative. Reducing the
baseline to a single BID representative (e.g.\ Randi\'c only)
would leave a high-dimensional BID residual space that any new
BID candidate could readily project onto, making the screen
trivially over-permissive on BID candidates. The verdict on the
non-BID alternative-family candidates of \S\ref{sec:novel_case}
(distance-, motif-, eccentricity-, centrality-, spectral-hybrid)
is unaffected by this composition choice, because those
candidates lie in functional subspaces that the BID-heavy
baseline cannot reach by construction; this is also why the
combined diagnostic admits non-zero passers on those families
(Table~\ref{tab:novel_combined}).

\begin{table}[h]
\centering
\caption{$30$-index classical baseline used for orthogonality screening.}
\label{tab:baseline}
\small
\begin{tabular}{ll}
\toprule
Family & Indices\\
\midrule
Degree-based (18) & $M_1$, $M_2$, $mM_1$, $mM_2$, $F$, $R$, SCI, $H$, $GA$, $AG$, ABC, ABS, AZI, $SO$, $SO_{\mathrm{red}}$, Alb, $\sigma$, $\mathrm{redM}_1$\\
Distance-based (8) & $W$, WW, HR, MTI (Schultz), Gutman, Mostar, $Sz$, $PI$\\
Spectral (3) & EE, Energy, Spectral Radius\\
Other (1) & Balaban $J$\\
\bottomrule
\end{tabular}
\end{table}

% =====================================================================
\section{The Edge-Degree-Pair Basis observation}\label{sec:bid}
% =====================================================================

We isolate a structural reason why the BID family on
hydrogen-suppressed molecular graphs is redundancy-prone. The
theorem is restated here for completeness because it is used
\emph{operationally} in the screening argument (in particular, as
the structural upper bound that licenses the redundancy
verdict on the Wuzi family in
Section~\ref{sec:wuzi_case}); the same theorem appears in the
companion graph-theory paper~\cite{Paper1_Wuzi}, where it sets
the structural ceiling for the Wuzi-specific mathematical
consequences. The novelty in the present paper is not the theorem
itself --- the linearity observation that underpins it is standard
in the mathematical-chemistry literature~\cite{BorovicaninDasFurtulaGutman2017Zagreb}
--- but its operational use inside the screening pipeline.

\begin{theorem}[Edge-Degree-Pair Basis]\label{thm:bid}
Let $\mathcal{G}_\Delta$ denote the class of finite simple graphs
with maximum degree at most $\Delta$. For any BID index
$I_f\in\bid$ and any $G\in\mathcal{G}_\Delta$,
\begin{equation}\label{eq:bid_linear}
I_f(G)=\sum_{1\le i\le j\le\Delta} f(i,j)\cdot m_{ij}(G),
\end{equation}
where $m_{ij}(G)=\bigl|\{uv\in E(G):\{d_u,d_v\}=\{i,j\}\}\bigr|$.
Restricted to $\mathcal{G}_4$, the vector space spanned by all BID
indices has dimension at most $10$.
\end{theorem}

\begin{proof}
We establish the linear-decomposition identity~\eqref{eq:bid_linear}
first and then deduce the dimension bound.

\emph{Step 1: linear decomposition.} Fix $G\in\mathcal{G}_\Delta$ and,
for each pair $(i,j)$ with $1\le i\le j\le\Delta$, define
$E_{ij}(G) = \{ uv\in E(G) \,:\, \{d_u,d_v\}=\{i,j\}\}$,
so that $m_{ij}(G)=|E_{ij}(G)|$. Because every edge has
$d_u,d_v\in\{1,\ldots,\Delta\}$, the family $\{E_{ij}(G)\}_{i\le j}$
partitions $E(G)$. Hence
\[
I_f(G) = \sum_{uv\in E(G)} f(d_u,d_v)
= \sum_{1\le i\le j\le\Delta}\;\sum_{uv\in E_{ij}(G)} f(d_u,d_v).
\]
For every $uv\in E_{ij}(G)$ the multiset $\{d_u,d_v\}$ equals
$\{i,j\}$ and $f$ is symmetric, so $f(d_u,d_v)=f(i,j)$ is constant
on the inner sum, which collapses to $f(i,j)\cdot m_{ij}(G)$ and
yields~\eqref{eq:bid_linear}.

\emph{Step 2: dimension bound.} For each admissible pair $(i,j)$
with $1\le i\le j\le\Delta$, regard the edge-count
$m_{ij}:\mathcal{G}_\Delta\to\mathbb{N}_{\ge 0}$ as a functional on
$\mathcal{G}_\Delta$. Equation~\eqref{eq:bid_linear} shows
$I_f = \sum_{i\le j} f(i,j)\cdot m_{ij}$ as functionals on
$\mathcal{G}_\Delta$, so
$\bid \subseteq \operatorname{span}_{\mathbb{R}}\{m_{ij}:1\le i\le j\le\Delta\}$.
A subspace of a span of $N_\Delta$ vectors has dimension at most
$N_\Delta$, hence
\[
\dim_{\mathbb{R}}\bid \le N_\Delta = \frac{\Delta(\Delta+1)}{2}
\]
by the standard multiset-counting identity. For $\Delta=4$ this is
$N_4=10$.
\end{proof}

\begin{remark}\label{rem:bid_credit}
The observation that BID indices are linear in $m_{ij}$ counts is
standard~\cite{BorovicaninDasFurtulaGutman2017Zagreb}, and the
finite-dimensional linear-algebraic view of vertex-degree-based
indices has recently been developed in general by
Rada~\cite{Rada2026LinearVDB}. The contribution here is
(i)~the explicit dimension-$10$ bound on $\mathcal{G}_4$ and
(ii)~its operational use in the screening pipeline: any new BID
index proposed for molecular chemistry must justify its
independence from the existing zoo within a $10$-dimensional
subspace.
\end{remark}

\begin{remark}\label{rem:rank_vs_bound}
The bound $\dim_{\mathbb{R}}\bid \le 10$ is a structural property
of the BID class on $\mathcal{G}_4$, not a property of any
specific dataset. Empirically, the observed PCA effective rank of
the $30$-index baseline on the four MoleculeNet datasets is much
smaller ($\mathrm{PC}_{95}=3$, $\mathrm{PC}_{99}=5$; see
Table~\ref{tab:baseline_redundancy}), reflecting the specific
distribution of molecular shapes in those datasets. The
theoretical bound is the ceiling that the empirical rank
cannot exceed under any dataset; the empirical rank is the
sharper, dataset-specific reality. Both numbers are useful: the
empirical rank quantifies redundancy on the dataset at hand,
while the theoretical bound tells the reader that no future
dataset of $\mathcal{G}_4$-molecules and no proposal of further
BID indices can raise the spanned dimension above ten.
\end{remark}

\begin{corollary}\label{cor:bid_redundancy}
On $\mathcal{G}_4$, any $11$ BID indices are linearly dependent.
Parametric BID families (general Randi\'c $R_\alpha$, general
sum-connectivity $\chi_\beta$, Wuzi, etc.) cannot escape the
$10$-dimensional ceiling by adding parameters.
\end{corollary}

In practice, the rank of the matrix $M$ with rows
$(m_{11}, m_{12}, \ldots, m_{44})$ over a chemistry dataset is
strictly less than $10$ because several $m_{ij}$ are sparse or
near-constant. For example, $m_{14}$ and $m_{44}$ require
quaternary carbons connected to methyl groups (rare),
$m_{11}$ requires an isolated bond (extremely rare in drug-like
chemistry).

% =====================================================================
\section{The screening protocol}\label{sec:protocol}
% =====================================================================

Let $X\in\mathbb{R}^{n\times p}$ be a matrix of $p$ baseline indices
on $n$ molecules, $y\in\mathbb{R}^n$ the target property, and
$z\in\mathbb{R}^n$ the candidate index values on the same molecules.
The screening pipeline produces four diagnostics.

\subsection{Pairwise correlation and redundancy threshold}
Compute the $p\times p$ Pearson correlation matrix $\mathrm{Cor}(X)$
and the vector of correlations $\mathrm{cor}(z,X_b)$ for
$b=1,\ldots,p$. The candidate is \emph{statistically redundant} if
$\max_b|\mathrm{cor}(z,X_b)|\ge\tau$ for a threshold $\tau$ (we
report $\tau=0.95$ as the primary threshold and also
$\tau=0.90$ as a secondary check). The redundant-pair count
inside the baseline,
$|\{(b_1,b_2):|\mathrm{Cor}(X)_{b_1,b_2}|\ge\tau\}|$, is a
companion summary statistic of baseline redundancy.

\subsection{Principal-component effective rank}
Compute the PCA~\cite{WoldEsbensenGeladi1987PCA} of the
standardized baseline $X$. Let
$\lambda_1\ge\lambda_2\ge\cdots\ge\lambda_p$ be the eigenvalues of
$\mathrm{Cor}(X)$ and $C_k=\sum_{i\le k}\lambda_i/\sum_i\lambda_i$
the cumulative variance fraction. The
$\mathrm{PC}_{1-\alpha}$ effective rank is the smallest $k$ with
$C_k\ge 1-\alpha$; we report $\mathrm{PC}_{95}, \mathrm{PC}_{99}$.

\subsection{Variance Inflation Factor}
For each baseline index $b$, $\mathrm{VIF}_b=(1-R_b^2)^{-1}$ where
$R_b^2$ is the coefficient of determination of $X_b$ regressed on
the other $p-1$ baseline columns. $\mathrm{VIF}_b>10$ is the
conventional ``high collinearity'' threshold~\cite{OBrien2007VIF};
$\mathrm{VIF}_b=\infty$ indicates exact linear dependence within
the baseline.

\subsection{Partial correlation with target}
Compute $z^{\mathrm{res}}=z-\hat z$ and $y^{\mathrm{res}}=y-\hat y$
where $\hat z, \hat y$ are the OLS fits regressing $z$ and $y$ on
the baseline $X$. The partial correlation
$\mathrm{pcor}(z,y\mid X)=\mathrm{cor}(z^{\mathrm{res}},y^{\mathrm{res}})$
measures the signal in $z$ relative to the target $y$ \emph{after}
controlling for everything the baseline already captures. We use
$|\mathrm{pcor}|<0.10$ as the conventional ``negligible
additional signal'' band. The cutoff $|\mathrm{pcor}|=0.10$
corresponds to $\mathrm{pcor}^2 = 0.01$, i.e.\ the candidate
explains an additional $1\%$ of \emph{residual} variance over
the baseline. For QSPR endpoints with bulk linear $R^2$ in the
$0.1$--$0.5$ range (cf.\ the LASSO benchmark of
Table~\ref{tab:matched_class}), a $1\%$ residual-$R^2$ pickup is
the smallest effect that is reliably distinguishable from noise
on samples of size $n \sim 10^3$. The threshold-sensitivity
table (Table~\ref{tab:threshold_sensitivity}, \S\ref{sub:threshold_sensitivity})
confirms that the qualitative verdict on the Wuzi family is
stable for tighter ($0.20$) and looser ($0.05$) cutoffs.

\subsection{Decision summary}

\begin{table}[h]
\centering
\caption{Screening diagnostic summary; satisfying both conditions is evidence of non-negligible residual linear signal, not a proof of QSPR utility.}
\label{tab:decision}
\begin{tabular}{lll}
\toprule
Diagnostic & Pass condition & Interpretation\\
\midrule
$\max_b|\mathrm{cor}(z,X_b)|$ & $<\tau=0.95$ & not strongly pairwise-redundant with any baseline index\\
$|\mathrm{pcor}(z,y\mid X)|$  & $\ge 0.10$  & non-negligible residual linear association with the target after controlling for the baseline\\
\bottomrule
\end{tabular}
\end{table}

The combined diagnostic of Table~\ref{tab:decision} should be read
as a \emph{conservative pre-flight diagnostic for independent linear
QSPR signal}, not as a hard feature-selection filter or a utility
certification. Passing the screen does not prove utility; failing
the screen is evidence of redundancy under the tested baseline,
dataset, and linear diagnostics. We deliberately avoid language
such as ``meaningful new index'' in favour of the more precise
``candidate with non-negligible residual linear signal under the
tested $30$-index baseline''. Threshold sensitivity is documented
in \S\ref{sub:threshold_sensitivity} and the Lipophilicity failure
mode --- where the strict combined diagnostic retains zero features
in $4/5$ ML-benchmark folds --- is treated as a formal limitation
in \S\ref{sub:limitations}.

% =====================================================================
\section{Results}\label{sec:results}
% =====================================================================

The two case studies and the downstream cross-validated ML
benchmark together establish that (i) the Wuzi parametric BID
family is fully redundant against the classical baseline on every
QSPR dataset we tested (Section~\ref{sec:wuzi_case}); (ii) the
pairwise $|r| < 0.95$ screen is over-permissive when applied to
candidates from non-BID design families
(Section~\ref{sec:novel_case}); and (iii) the pairwise screen
preserves predictive performance within cross-validation noise on
four MoleculeNet benchmarks (Section~\ref{sec:ml_benchmark}). All
three results are reproducible from the public repository, and
the full per-dataset numerical tables are in
\texttt{results/}.

\subsection{Case study 1: the Wuzi parametric family}\label{sec:wuzi_case}

The Wuzi index family is the three-parameter BID family
$W(G;\alpha,\beta,\gamma) = \sum_{uv\in E(G)} (d_u d_v)^\alpha (d_u + d_v)^\beta
\exp(\gamma\cdot|d_u-d_v|/(d_u+d_v))$ developed in the companion
paper~\cite{Paper1_Wuzi}; we treat it here as one parametric BID
candidate and screen it on a $100$-point grid
$(\alpha, \beta, \gamma) \in \{-1, -1/2, 0, 1/2, 1\}^2 \times \{0, 1/2, 1, 2\}$
on each of the three regression QSPR datasets.

\subsubsection{Baseline redundancy on the three datasets}

\begin{table}[h]
\centering
\small
\caption{Baseline redundancy across the three QSPR datasets. Pair counts are out of $435$ (the number of unordered pairs of $30$ indices). $\mathrm{PC}_{\alpha}$ is the smallest number of principal components whose cumulative explained variance is at least $\alpha\%$ of the baseline.}
\label{tab:baseline_redundancy}
\begin{tabular}{lrrrrrr}
\toprule
Dataset & $n$ & $\lvert r\rvert \ge 0.90$ & $\lvert r\rvert \ge 0.95$ & $\mathrm{PC}_{90}$ & $\mathrm{PC}_{95}$ & $\mathrm{PC}_{99}$\\
\midrule
ESOL          & $1127$ & $203$ ($46.7\%$) & $159$ ($36.6\%$) & $2$ & $3$ & $5$\\
FreeSolv      & $\phantom{0}639$ & $257$ ($59.1\%$) & $162$ ($37.2\%$) & $2$ & $3$ & $5$\\
Lipophilicity & $4200$ & $200$ ($46.0\%$) & $164$ ($37.7\%$) & $2$ & $3$ & $5$\\
\bottomrule
\end{tabular}
\end{table}

Across three independent datasets the pattern is consistent: about
$37\%$ of the $435$ baseline pairs are above the $0.95$ threshold,
and three principal components capture $95\%$ of the variance
(Figs.~\ref{fig:pca},~\ref{fig:redundancy_bars}).
Figure~\ref{fig:corr_heatmap} shows the full
$30 \times 30$ baseline correlation matrix on ESOL; the dense red
block in the upper-left corner is the degree-based index family,
which drives most of the redundancy reported in
Table~\ref{tab:baseline_redundancy}.

\begin{figure}[h]
\centering
\includegraphics[width=0.75\linewidth]{fig2_pca_scree.pdf}
\caption{Cumulative explained variance of PCA on the $30$-index baseline.}
\label{fig:pca}
\end{figure}

\begin{figure}[h]
\centering
\includegraphics[width=0.65\linewidth]{fig3_redundancy_bars.pdf}
\caption{Redundant index-pair counts out of $435$ baseline pairs (unordered pairs of $30$ indices) at two thresholds.}
\label{fig:redundancy_bars}
\end{figure}

\begin{figure}[h]
\centering
\includegraphics[width=0.6\linewidth]{fig8_correlation_heatmap.pdf}
\caption{$30\times 30$ baseline correlation matrix on ESOL. The dense red core is the degree-based index block.}
\label{fig:corr_heatmap}
\end{figure}

\subsubsection{Wuzi-family screening verdict and the multi-dimensional redundancy at $(-1,-1,1)$}

We screen the Wuzi family on $100$ grid points per dataset
($300$ grid points across the three QSPR datasets) and report the
compact per-dataset verdict in Table~\ref{tab:wuzi_verdict_paper1}.
We isolate the parameter triple
$(\alpha, \beta, \gamma) = (-1, -1, 1)$ below because it is the
single methodologically central grid point of the entire sweep
and illustrates the necessary-but-not-sufficient nature of the
pairwise screen.

\begin{table}[h]
\centering
\small
\caption{Wuzi $100$-point grid screening verdict across three QSPR datasets. ``Pass'' counts grid points at which $\max_{b}\lvert r(W, X_b)\rvert$ is strictly below $0.95$ for every classical baseline index $X_b$. ``Min $\max\lvert r\rvert$'' is the across-grid minimum of $\max_b\lvert r\rvert$ (the best-passing grid point). ``Max $\lvert\mathrm{pcor}\rvert$'' is the across-grid maximum of $\lvert\mathrm{pcor}(W, y \mid X)\rvert$, the partial correlation of $W$ with the QSPR target $y$ after regressing $W$ on the $30$-index baseline $X$.}
\label{tab:wuzi_verdict_paper1}
\begin{tabular}{lcccc}
\toprule
Dataset & Pass (max $\lvert r\rvert < 0.95$) & Fail & Min max $\lvert r\rvert$ & Max $\lvert\mathrm{pcor}(W, y \mid X)\rvert$\\
\midrule
ESOL          & 0 & 100 & $0.965$ & $0.028$\\
FreeSolv      & 1 &  99 & $0.943$ & $0.036$\\
Lipophilicity & 0 & 100 & $0.965$ & $0.009$\\
\bottomrule
\end{tabular}
\end{table}

\textbf{Verdict.} Across the $300$ grid points the Wuzi family
clears the pairwise $\max|r| < 0.95$ screen at exactly one grid
point (FreeSolv at $(-1, -1, 1)$, $\max|r| = 0.943$ with $mM_2$).
The maximum partial correlation with the target across the $300$
grid points is $0.036$ (FreeSolv); the corresponding maxima on
ESOL and Lipophilicity are $0.028$ and $0.009$ respectively, all
below $0.10$. By the combined diagnostic of
Table~\ref{tab:decision}, no Wuzi parameter triple in the grid
passes the diagnostic for non-negligible residual linear signal
under the tested $30$-index baseline.

\textbf{The multi-dimensional redundancy at $(-1, -1, 1)$.}
At this parameter triple the Wuzi value is exactly captured by a
linear combination of the $30$ baseline indices on \emph{all three}
QSPR datasets (the regression residual has variance below numerical
precision on ESOL, FreeSolv, and Lipophilicity alike), so the
partial correlation with the respective target is identically zero
in each case. The pairwise $\max|r|$ at this parameter triple is
dataset-dependent ($0.965, 0.943, 0.977$ for ESOL / FreeSolv / Lipo;
$mM_2$ is the most-correlated baseline in all three cases): only
the FreeSolv value falls below the conventional pairwise threshold,
but the multi-dimensional redundancy is structural and identical
across datasets. This illustrates the methodological point of
Section~\ref{sec:protocol}: the pairwise $|r|$ screen catches only
one-dimensional redundancy, and a parameter setting that ``passes''
the pairwise screen because of one dataset's particular degree
distribution can still be exactly captured by the baseline in the
regression sense. Figure~\ref{fig:wuzi_grid_all} shows the full
$\max|r|$ grid across all three datasets.

\begin{figure}[h]
\centering
\includegraphics[width=\linewidth]{fig5b_wuzi_param_heatmaps_all.pdf}
\caption{$\max|r|$ between the Wuzi index at each $(\alpha,\beta,\gamma)$
grid point and any of the $30$ classical baseline indices, across
the three datasets (rows) and four $\gamma$ slices (columns).
The starred cell (FreeSolv, $(-1,-1,1)$) is the only grid point
that falls below $|r|=0.95$; at the same parameter triple the
partial correlation with the target is identically zero on all three
datasets.}
\label{fig:wuzi_grid_all}
\end{figure}

\subsubsection{Consistency with the BID-basis observation}

Corollary~\ref{cor:bid_redundancy} predicts that no parametric BID
family can have effective dimension above $10$ on $\mathcal{G}_4$.
With three principal components capturing $95\%$ of the baseline
variance, the empirical effective rank is much lower than $10$; a
new BID index has at most a $7$-dimensional ``budget''
within which to be statistically independent of the leading PCA
factors. The empirical screening verdict that no Wuzi parameter
point clears $|r|<0.95$ \emph{and} retains
$|\mathrm{pcor}|\ge 0.10$ with the QSPR target is consistent
with this structural constraint.

\subsubsection{Threshold sensitivity of the Wuzi verdict}\label{sub:threshold_sensitivity}

A standard concern with any thresholded screen is whether the
qualitative verdict survives reasonable variation in the cutoffs
$(\tau_r, \tau_{\mathrm{pcor}})$. Table~\ref{tab:threshold_sensitivity}
tabulates the Wuzi $100$-point grid pairwise-pass counts at
$\tau_r \in \{0.90, 0.95, 0.99\}$ and the combined-diagnostic-pass
counts across the full $3 \times 3$ grid of
$\tau_r \in \{0.90, 0.95, 0.99\}$ and
$\tau_{\mathrm{pcor}} \in \{0.05, 0.10, 0.20\}$
on each of the three QSPR datasets (full table in
\texttt{results/paper2\_threshold\_sensitivity.csv}).

\begin{table}[!htbp]
\centering
\small
\caption{Wuzi $100$-point grid: pass counts under threshold variation. Top block: pairwise screen only ($\max_b\lvert r(W, X_b)\rvert < \tau_r$). Bottom block: combined diagnostic (pairwise pass \emph{and} $\lvert\mathrm{pcor}(W, y \mid X)\rvert \ge \tau_{\mathrm{pcor}}$) across the full $3 \times 3$ threshold grid.}
\label{tab:threshold_sensitivity}
\begin{tabular}{lccc}
\toprule
& \multicolumn{3}{c}{Pairwise pass count ($\tau_r$)}\\
\cmidrule(lr){2-4}
Dataset & $\tau_r = 0.90$ & $\tau_r = 0.95$ & $\tau_r = 0.99$\\
\midrule
ESOL          & $0$ & $0$ & $27$\\
FreeSolv      & $0$ & $1$ & $37$\\
Lipophilicity & $0$ & $0$ & $30$\\
\bottomrule
\end{tabular}

\vspace{0.6em}

\begin{tabular}{l|ccc|ccc|ccc}
\toprule
& \multicolumn{3}{c|}{$\tau_r = 0.90$} & \multicolumn{3}{c|}{$\tau_r = 0.95$} & \multicolumn{3}{c}{$\tau_r = 0.99$}\\
$\tau_{\mathrm{pcor}}$ & $.05$ & $.10$ & $.20$ & $.05$ & $.10$ & $.20$ & $.05$ & $.10$ & $.20$\\
\midrule
ESOL          & $0$ & $0$ & $0$ & $0$ & $0$ & $0$ & $0$ & $0$ & $0$\\
FreeSolv      & $0$ & $0$ & $0$ & $0$ & $0$ & $0$ & $0$ & $0$ & $0$\\
Lipophilicity & $0$ & $0$ & $0$ & $0$ & $0$ & $0$ & $0$ & $0$ & $0$\\
\bottomrule
\end{tabular}

\vspace{0.3em}
\noindent\small{Combined-diagnostic pass count of the Wuzi $100$-point grid across the full $\{0.90, 0.95, 0.99\} \times \{0.05, 0.10, 0.20\}$ threshold grid (nine $(\tau_r, \tau_{\mathrm{pcor}})$ combinations).}
\end{table}
\FloatBarrier

\noindent The qualitative conclusion is stable across the tested
thresholds: \emph{no} Wuzi grid point clears the combined
diagnostic at any of the nine $(\tau_r, \tau_{\mathrm{pcor}})$
combinations in $\{0.90, 0.95, 0.99\} \times \{0.05, 0.10, 0.20\}$
on any of the three datasets. Relaxing the pairwise threshold to
$\tau_r = 0.99$ admits roughly a third of grid points through the
pairwise screen ($27$, $37$, $30$ of $100$), but none of those
survive even the loosest partial-correlation cutoff
$\tau_{\mathrm{pcor}} = 0.05$. The exact pass counts are
threshold-dependent; the conclusion that the Wuzi family lacks
independent linear QSPR signal on these datasets under the
$30$-index baseline is not.

\subsection{Case study 2: alternative-family candidates across four datasets}\label{sec:novel_case}

\noindent\emph{Scope note.} This case study screens twenty
alternative-family candidate indices on all four benchmarks
(ESOL, FreeSolv, Lipophilicity, BBBP), so the headline pattern
``$15$--$17/20$ clear the pairwise screen but only $0$--$4/20$
also clear the partial-correlation screen'' is a cross-dataset
finding, not a single-dataset case study. The candidates are
identical across datasets; the per-dataset verdict differs because
the residual signal under the $30$-index baseline differs by
target. We make no claim that any individual passer is a useful
descriptor; the screening verdict is a redundancy diagnostic, not
a chemical-utility certification. The candidates were chosen as
five distinct design philosophies (information-theoretic,
centrality, motif, eccentricity, spectral-hybrid) to test that
the screening conclusion generalises beyond the BID-family case
study of \S\ref{sec:wuzi_case}.

We screen $20$ candidate indices on the four datasets, drawn from
the five alternative design families listed in
Table~\ref{tab:candidates} and described below.

\begin{table}[h]
\centering
\caption{Twenty candidate indices grouped by design family.}
\label{tab:candidates}
\begin{tabular}{ll}
\toprule
Family & Candidate indices (working names)\\
\midrule
Information-theoretic (5) & InfoH\_deg, InfoH\_dist, InfoH\_spec, InfoH\_ecc, Bonchev\_Id\\
Centrality (5) & Btw\_sum, Cls\_sum, Eig\_sum, Harm\_cent, Btw\_var\\
Motif (4) & Triangles, MeanClust, FourCyc, CutVerts\\
Eccentricity (3) & TotEcc, Radius, MeanEcc\\
Spectral-hybrid (3) & AlgConn, Rho$\times$Dia, EE$\times$W\\
\bottomrule
\end{tabular}
\end{table}

\noindent\emph{Disclaimer.} None of the $20$ candidates is claimed
to be newly invented. Each has prior chemistry literature (the
information-theoretic group draws on Cao--Dehmer--Shi and related
entropy-style proposals; the centrality group on textbook graph
centrality measures; the motif and eccentricity groups on
standard graph invariants; the spectral-hybrid group on
combinations of standard spectral quantities). They appear here
purely as screening targets to test whether candidates from
different design philosophies escape the redundancy ceiling on
ESOL; the full per-candidate citation list is collected in
\texttt{docs/literature\_notes.md} in the companion repository.

\subsubsection{Screening verdict across the four datasets}

Table~\ref{tab:novel_combined} reports the per-dataset pairwise and
combined pass counts. The pairwise $|r|<0.95$ screen alone admits
$15$--$17$ of the $20$ candidates on every dataset; the combined
diagnostic that additionally requires
$|\mathrm{pcor}|\ge 0.10$ admits only $0$--$4$. The pattern is
consistent: the pairwise screen is uniformly over-permissive, and
the combined diagnostic is sharper on FreeSolv ($1/20$ pass),
BBBP ($2/20$), and Lipophilicity ($0/20$) than on ESOL ($4/20$).

For BBBP (binary target) the partial-correlation diagnostic is
computed via the linear probability model: we fit
$y \sim X_{\mathrm{baseline}}$ as ordinary least squares against
the $0/1$ class indicator and take residuals; the candidate is
then residualised by the same baseline and Pearson-correlated
with the residual class indicator. This is the standard
LPM-based partial correlation; it puts the BBBP row of
Table~\ref{tab:novel_combined} on the same scale as the three
regression rows. A logistic-pseudo-$R^2$ alternative
(Tjur's coefficient) is recorded as future work in
\S\ref{sub:limitations}.

\begin{table}[h]
\centering
\caption{Combined screening verdict on the $20$ candidates against the $30$-index baseline on each of the four MoleculeNet benchmarks. $\mathrm{PASS}_{\mathrm{pair}}$ = max $|r| < 0.95$; $\mathrm{PASS}_{\mathrm{combined}}$ = max $|r| < 0.95$ AND $|\mathrm{pcor}|\ge 0.10$. Per-dataset combined passers are listed in the footnote.}
\label{tab:novel_combined}
\begin{tabular}{lcccr}
\toprule
Dataset & $n$ & $\mathrm{PASS}_{\mathrm{pair}}$ / $20$ & $\mathrm{PASS}_{\mathrm{combined}}$ / $20$ & $\max|\mathrm{pcor}|$\\
\midrule
ESOL          & $1\,127$ & $15$ & $4$ & $0.154$\\
FreeSolv      & $\phantom{0}639$  & $15$ & $1$ & $0.129$\\
Lipophilicity & $4\,200$ & $16$ & $0$ & $0.077$\\
BBBP$^\ddag$  & $2\,039$ & $17$ & $2$ & $0.225$\\
\bottomrule
\end{tabular}

\vspace{0.4em}
\noindent\small{Combined passers per dataset:
ESOL: $\mathrm{InfoH\_deg}\,(0.154)$, CutVerts $(0.118)$, FourCyc $(0.112)$, MeanEcc $(0.109)$.
FreeSolv: Triangles $(0.107)$.
Lipophilicity: none.
BBBP: $\mathrm{InfoH\_deg}\,(0.225)$, $\mathrm{InfoH\_dist}\,(\approx 0.16)$.
$^\ddag$ BBBP partial correlation computed via LPM (see text).}
\end{table}

The headline pattern: the pairwise $|r|<0.95$ screen alone is
over-permissive on every dataset --- $75$--$85\%$ of candidates
clear it --- but the combined diagnostic that additionally
requires $|\mathrm{pcor}|\ge 0.10$ leaves $0$--$4$ candidates as
plausible carriers of independent linear signal under the
$30$-index baseline. The strongest passer across the four
datasets is $\mathrm{InfoH\_deg}$ on BBBP at $|\mathrm{pcor}|=0.225$;
no other candidate reaches $|\mathrm{pcor}|>0.2$ on any dataset.
Detailed per-candidate numbers are in
\texttt{results/novel\_candidates\_experiment/novel\_candidates\_multidataset.csv}.

These per-dataset combined passers should be read as
``residual-signal positive on that specific target under the
$30$-index linear baseline'', not as useful QSPR descriptors;
the diagnostic is a redundancy filter, not a utility
certification.

\subsubsection{Structural interpretation of the four ESOL passers}\label{sec:passer_analysis}

We focus the structural reading on ESOL because the four ESOL
passers span the widest set of design philosophies among the four
datasets and so best illustrate \emph{why} the partial-correlation
diagnostic catches signal that pairwise correlation misses. The
qualitative argument transfers to the FreeSolv (Triangles) and BBBP
($\mathrm{InfoH\_deg}$, $\mathrm{InfoH\_dist}$) passers; the
specific structural reading below is ESOL-only.

The four candidates that clear the combined diagnostic on ESOL span
three distinct design philosophies, and each computes a graph
property that is \emph{not} expressible as a sum of per-edge
contributions:

\begin{itemize}
\item $\mathrm{InfoH\_deg}$, the Shannon entropy of the degree
multiplicities of $G$, is a logarithmic functional of the degree
\emph{distribution}, not of the degree sequence summed edgewise.
\item $\mathrm{CutVerts}$, the number of articulation points of $G$,
is a vertex-connectivity invariant: the count of vertices whose
removal disconnects $G$. It depends on the global block structure
of $G$, not on local edge properties.
\item $\mathrm{FourCyc}$, the number of induced $4$-cycles, captures
a cycle-counting property of $G$. On hydrogen-suppressed molecular
graphs of order $n \le 30$ it is small (often zero) and captures
ring-substructure information disjoint from degree-based measures.
\item $\mathrm{MeanEcc}$, the mean of the vertex eccentricities, is
a distance-based scalar that depends on shortest-path lengths
rather than on degrees.
\end{itemize}

This shared structural feature explains why all four can clear the
partial-correlation threshold even when their pairwise correlation
against an individual baseline is high. $\mathrm{MeanEcc}$, for
example, has $\max |r| = 0.930$ against the modified second
Zagreb index $mM_2$ on ESOL --- well above any natural redundancy
threshold for a pair of degree-based indices --- but it still
retains $|\mathrm{pcor}| = 0.109$ once the full $30$-index linear
span is regressed out. Pairwise correlation against a single
baseline does not preclude residual signal orthogonal to the
$30$-dimensional baseline subspace, and the partial-correlation
diagnostic is the component of the combined diagnostic that catches
this. By contrast, the $11$ candidates that clear the pairwise
screen with $|\mathrm{pcor}| < 0.10$ are well-approximated by a
linear combination of the $30$ baseline indices: high pairwise
$|r|$ against \emph{one} baseline is neither necessary nor
sufficient for redundancy against the \emph{full} baseline.

We emphasise that the four passing partial correlations are small
in absolute terms ($0.109 \le |\mathrm{pcor}| \le 0.154$, all under
$0.2$). The reading we recommend is structural: the partial-correlation
diagnostic is \emph{consistent with} the hypothesis that these
candidates capture target-relevant information not linearly
explained by the $30$-index baseline. Strictly proving subspace
orthogonality would require an explicit basis decomposition we do
not undertake here; what the data establish is the weaker statement
that residual linear association with the target persists after
projecting out the baseline. Whether any of
$\mathrm{InfoH\_deg}$, $\mathrm{CutVerts}$, $\mathrm{FourCyc}$,
or $\mathrm{MeanEcc}$ actually improves QSPR performance on $\log S$
or any other endpoint is a separate question, answered in part by
the ML benchmark of Section~\ref{sec:ml_benchmark} and ultimately
requiring downstream validation on tasks beyond ESOL.

\subsection{Downstream ML benchmark: does screening preserve predictive performance?}\label{sec:ml_benchmark}

The case studies of Sections~\ref{sec:wuzi_case} and~\ref{sec:novel_case}
demonstrate that the combined diagnostic is more selective than the
pairwise screen alone. They do not yet demonstrate that the criterion
is useful for building predictive QSPR models. This section closes
that gap: we train a RandomForest (300 trees, fixed seed,
\texttt{class\_weight="balanced"} for classification) on each of four
datasets under four feature configurations across $5$-fold
cross-validation, and report whether the screened feature set
matches the full-baseline model's predictive performance within
CV noise.

\textbf{Datasets.} The three QSPR regression benchmarks of
Section~\ref{sec:datasets} are augmented with a fourth dataset to
test that the pipeline generalizes beyond regression:
\textbf{BBBP} (\cite{Martins2012BBBP} via~\cite{MoleculeNet}), $2\,050$
molecules in the source CSV of which $2\,039$ parse via RDKit and
enter the screen; binary target ``blood-brain barrier
penetration''; class distribution approximately $3{:}1$
(positive:negative), handled with a class-balanced RandomForest.

\textbf{Configurations and cross-validation protocol.} $5$-fold
cross-validation, stratified for the classification task, identical
folds across all configurations so the comparison isolates the
effect of feature selection. All feature selection is performed
\emph{inside each training fold} to avoid test-fold information
leakage.
\begin{itemize}
\item \textbf{full}: all $30$ baseline indices.
\item \textbf{pca\_95}: top-$k$ principal components of the
training-fold-standardized baseline at $95\%$ cumulative variance
($k=3$ on all four datasets, every fold).
\item \textbf{pairwise\_pruned}: greedy redundancy filter at
$|r| < 0.95$ on the training fold (keep one representative per
correlation cluster).
\item \textbf{combined\_pruned}: pairwise filter plus
$|\mathrm{pcor}(z, y \mid X_{\mathrm{others}})| \ge 0.10$ on the
training fold's pairwise-pruned set.
\end{itemize}

\textbf{Results.} Table~\ref{tab:ml_benchmark} and
Figure~\ref{fig:ml_benchmark} summarize the per-dataset, per-configuration
RandomForest performance (mean $\pm$ std across the $5$ folds) and
mean feature count.

\begin{table}[h]
\centering
\small
\caption{RandomForest $5$-fold cross-validation performance under four feature configurations. Regression metric is RMSE (lower is better); classification metric is ROC-AUC (higher is better). Reported as mean $\pm$ std across folds. The \texttt{combined\_pruned} screen retained zero features in $4$ of $5$ folds on Lipophilicity (the RMSE shown is from the single fold that retained one feature).}
\label{tab:ml_benchmark}
\begin{tabular}{lccccrr}
\toprule
Dataset & \texttt{full} & \texttt{pca\_95} & \texttt{pairwise\_pruned} & \texttt{combined\_pruned} & full $n$ & pairwise $n$\\
\midrule
ESOL          & $1.437 \pm 0.102$ & $1.455 \pm 0.096$ & $1.435 \pm 0.087$ & $1.475 \pm 0.123$ & $30$ & $7.4$ ($7$--$9$)\\
FreeSolv      & $3.537 \pm 0.228$ & $3.544 \pm 0.264$ & $3.540 \pm 0.224$ & $3.561 \pm 0.300$ & $30$ & $8.0$\\
Lipophilicity & $1.030 \pm 0.024$ & $1.065 \pm 0.018$ & $1.028 \pm 0.022$ & $1.287$ ($1/5$ folds) & $30$ & $8.0$\\
BBBP          & $0.846 \pm 0.006$ & $0.818 \pm 0.014$ & $0.860 \pm 0.010$ & $0.840 \pm 0.010$ & $30$ & $6.0$\\
\bottomrule
\end{tabular}
\end{table}

\begin{figure}[h]
\centering
\includegraphics[width=\linewidth]{fig9_ml_benchmark.pdf}
\caption{Downstream ML benchmark on four datasets, $5$-fold cross
validation. Top row: primary metric (RMSE for the three regression
panels, ROC-AUC for BBBP), mean $\pm$ std across folds. Bottom row:
mean feature count per configuration, error bar = min--max range
across folds. The pairwise screen matches the full-baseline
RandomForest on every dataset, with mean differences below CV
noise, while using roughly $4\times$ (range $3.75$--$5\times$)
fewer features. On Lipophilicity the strict combined screen retains
zero features in $4$ of $5$ folds.}
\label{fig:ml_benchmark}
\end{figure}

\textbf{Interpretation.} On the four MoleculeNet benchmarks tested,
the \texttt{pairwise\_pruned} configuration matches \texttt{full}
within CV noise: mean RMSE differences are $|\Delta| \le 0.005$ on
all three regression datasets and the mean and std intervals
overlap. On BBBP, \texttt{pairwise\_pruned} attains ROC-AUC
$0.860 \pm 0.010$ versus \texttt{full} at $0.846 \pm 0.006$; the
mean delta is $0.014$, with marginal overlap of the std bands.
This is best read as ``matches or marginally exceeds within CV
noise'' rather than a definitive improvement. The PCA-95
configuration consistently underperforms by $\approx 0.02$ RMSE or
$\approx 0.03$ ROC-AUC, suggesting that the principal directions
of the standardized $30$-index baseline are not the
target-relevant directions for these specific QSPR problems.

The stricter \texttt{combined\_pruned} configuration is more
aggressive and on the regression datasets shows a small performance
penalty (mean RMSE up to $0.04$ higher than \texttt{full}, still
overlapping std bands). On Lipophilicity the combined screen
retains zero features in $4$ of the $5$ folds, an honest failure
mode that reflects the dataset's weak topology-to-$\log D$
partial-correlation structure (no pairwise-pruned index has
$|\mathrm{pcor}| \ge 0.10$ with $\log D$ on most folds). The
threshold $|\mathrm{pcor}| = 0.10$ should be relaxed (or replaced
with a target-correlation rank cut) on datasets where the underlying
topology-target relationship is weak. We do not recommend the
combined screen as a hard filter; it is best read as a diagnostic
that flags ``low pairwise correlation but low partial correlation
too'' --- which is exactly the failure mode the Wuzi family at
$(-1,-1,1)$ exhibits.

\subsubsection{Comparison to LASSO, ElasticNet, mRMR, Boruta, and L1-logistic feature selection}\label{sub:lasso_comparator}

To position the pairwise / combined diagnostic against a broader
set of feature-selection comparators, we evaluate five external
methods alongside the pipeline configurations of
Table~\ref{tab:ml_benchmark}:
\begin{itemize}
\item LASSO with $5$-fold CV-tuned $\alpha$ (three regression datasets);
\item ElasticNet with CV-tuned $\alpha$ and $\ell_1$-ratio
$\in \{0.3, 0.5, 0.7, 0.9\}$ (three regression datasets);
\item mRMR top-$k$ with $k$ matched to the pairwise-pruned mean
feature count of Table~\ref{tab:ml_benchmark} (all four datasets);
\item Boruta with RandomForest base estimator (all four datasets);
\item $\ell_1$-Logistic with CV-tuned $C$ (BBBP only; the
classification analog of LASSO).
\end{itemize}
The cached $30$-index baseline matrix and the same $5$-fold splits
as Table~\ref{tab:ml_benchmark} are used; for mRMR and Boruta we
evaluate the selected subset under ordinary linear regression
(regression) or logistic regression (BBBP) to keep the comparator
linear-model class consistent with LASSO. The screen-pipeline rows
remain RandomForest, so the comparison is intentionally across two
model classes; we read it as a sanity check that the pairwise screen
is not being compared to a strawman, not as a head-to-head model
benchmark.

\begin{table}[h]
\centering
\small
\caption{Screen-pipeline configurations (RandomForest, rows from Table~\ref{tab:ml_benchmark}) vs.\ five external feature-selection comparators (LASSO, ElasticNet, mRMR, Boruta, $\ell_1$-Logistic) on the four MoleculeNet benchmarks. Regression metric is RMSE (lower is better); classification metric for BBBP is ROC-AUC (higher is better). Reported as mean across 5 folds with the same seed.}
\label{tab:fs_comparison}
\begin{tabular}{llrr}
\toprule
Dataset & Configuration & Metric & Mean $n$ features\\
\midrule
ESOL & \texttt{full} (RF)                 & RMSE $1.437$ & $30.0$\\
     & \texttt{pairwise\_pruned} (RF)     & RMSE $1.435$ & $\phantom{0}7.4$\\
     & \texttt{combined\_pruned} (RF)     & RMSE $1.475$ & $\phantom{0}4.4$\\
     & LASSO (linear, CV)                 & RMSE $1.545$ & $14.6$\\
     & ElasticNet (linear, CV)            & RMSE $1.546$ & $16.2$\\
     & mRMR $k{=}7$ (linear)              & RMSE $1.654$ & $\phantom{0}7.0$\\
     & Boruta + RF (linear scoring)       & RMSE $1.620$ & $10.0$\\
     & \emph{pairwise} $\to$ \emph{LASSO} (matched$^{\ddagger}$) & RMSE $1.580$ & $\phantom{0}7.2$\\
\midrule
FreeSolv & \texttt{full} (RF)             & RMSE $3.537$ & $30.0$\\
         & \texttt{pairwise\_pruned} (RF) & RMSE $3.540$ & $\phantom{0}8.0$\\
         & \texttt{combined\_pruned} (RF) & RMSE $3.561$ & $\phantom{0}3.2$\\
         & LASSO (linear, CV)             & RMSE $3.527$ & $13.4$\\
         & ElasticNet (linear, CV)        & RMSE $3.515$ & $18.6$\\
         & mRMR $k{=}8$ (linear)          & RMSE $3.618$ & $\phantom{0}8.0$\\
         & Boruta + RF (linear scoring)   & RMSE $3.588$ & $\phantom{0}2.0$\\
         & \emph{pairwise} $\to$ \emph{LASSO} (matched$^{\ddagger}$) & RMSE $3.465$ & $\phantom{0}6.4$\\
\midrule
Lipophilicity & \texttt{full} (RF)        & RMSE $1.030$ & $30.0$\\
              & \texttt{pairwise\_pruned} (RF) & RMSE $1.028$ & $\phantom{0}8.0$\\
              & \texttt{combined\_pruned} (RF) & RMSE $1.287^{\,\dagger}$ & $\phantom{0}0.2$\\
              & LASSO (linear, CV)        & RMSE $1.130$ & $15.6$\\
              & ElasticNet (linear, CV)   & RMSE $1.130$ & $17.8$\\
              & mRMR $k{=}8$ (linear)     & RMSE $1.142$ & $\phantom{0}8.0$\\
              & Boruta + RF (linear scoring) & RMSE $1.144$ & $\phantom{0}5.0$\\
              & \emph{pairwise} $\to$ \emph{LASSO} (matched$^{\ddagger}$) & RMSE $1.137$ & $\phantom{0}7.0$\\
\midrule
BBBP & \texttt{full} (RF)                 & AUC $0.846$ & $30.0$\\
     & \texttt{pairwise\_pruned} (RF)     & AUC $0.860$ & $\phantom{0}6.0$\\
     & \texttt{combined\_pruned} (RF)     & AUC $0.840$ & $\phantom{0}4.0$\\
     & $\ell_1$-Logistic (CV)             & AUC $0.790$ & $13.6$\\
     & mRMR $k{=}6$ (logistic)            & AUC $0.747$ & $\phantom{0}6.0$\\
     & Boruta + RF (logistic scoring)     & AUC $0.791$ & $23.0$\\
\bottomrule
\end{tabular}

\vspace{0.5em}
\noindent\small{$^\dagger$ Strict combined screen degenerates: retains zero features in $4/5$ folds; the RMSE shown is from the single fold that retained one feature. See \S\ref{sub:limitations}.}
\noindent\small{$^\ddagger$ Pipeline-then-LASSO: in-fold pairwise screen $\to$ LASSO on the surviving features, same model class throughout. See \S\ref{sub:pipeline_then_lasso} and Table~\ref{tab:matched_class}.}
\end{table}

The headline reading is qualified by the cross-model-class
nature of the comparison. The screen-pipeline rows use a
RandomForest, and the comparators use linear regression
(regression datasets) or logistic regression (BBBP); the
\texttt{pairwise\_pruned} RandomForest performs competitively
with the tested external feature-selection comparators under this
mixed-model comparison, with feature counts in the range $6$--$8$
versus $6$--$23$ for the comparators. On ESOL,
\texttt{pairwise\_pruned} has lower RMSE than LASSO in this
benchmark ($1.435$ at $7.4$ features versus $1.545$ at $14.6$
features) and lower RMSE than mRMR-$k{=}7$ ($1.654$); on BBBP,
\texttt{pairwise\_pruned} has higher ROC-AUC than $\ell_1$-Logistic
($0.860$ at $6$ features versus $0.790$ at $13.6$ features) and
higher AUC than mRMR-$k{=}6$ ($0.747$). Because the comparators (LASSO, ElasticNet, mRMR, Boruta,
$\ell_1$-Logistic) and the screen-pipeline rows
(\texttt{full}, \texttt{pairwise\_pruned}, \texttt{combined\_pruned})
use different model classes (linear vs.\ RandomForest), the
single most honest reading of this table is the same-model-class
``pipeline-then-LASSO'' rows reported in
\S\ref{sub:pipeline_then_lasso} below. Full per-fold standard
deviations for all comparators are in
\texttt{results/paper2\_feature\_selection\_comparison.csv}.

\subsubsection{Matched-model-class benchmark: pipeline-then-LASSO}\label{sub:pipeline_then_lasso}

To eliminate the cross-model-class qualification on
Table~\ref{tab:fs_comparison} we re-run the screen-pipeline /
no-screen-pipeline comparison entirely within the LASSO
(regression) and $\ell_1$-Logistic (classification) model classes.
For each of the four datasets, under the same fold seed and the
same $5$-fold splits as Table~\ref{tab:ml_benchmark} and
Table~\ref{tab:fs_comparison}, we evaluate three configurations:

\begin{itemize}
\item \texttt{full (LASSO/L1-Log)}: LASSO (or $\ell_1$-Logistic on
BBBP) with $5$-fold-CV-tuned regularisation, fit on the full
$30$-index baseline inside each training fold.
\item \texttt{pairwise} $\to$ \texttt{LASSO/L1-Log}: pairwise
greedy redundancy filter at $|r|<0.95$ applied inside each
training fold, then LASSO (or $\ell_1$-Logistic) fit on the
surviving features.
\item \texttt{combined} $\to$ \texttt{LASSO/L1-Log}: pairwise
filter plus the in-fold partial-correlation filter
$|\mathrm{pcor}|\ge 0.10$, then LASSO (or $\ell_1$-Logistic) fit on
the survivors.
\end{itemize}

The pairwise / combined selectors are the byte-identical functions
used by the RandomForest benchmark of script
\texttt{11\_ml\_benchmark.py}, so the surviving feature subsets are
exactly the ones the RF benchmark sees. The only difference is the
downstream model class: LASSO / $\ell_1$-Logistic instead of
RandomForest.

Table~\ref{tab:fs_comparison} includes a single compact
\emph{pairwise} $\to$ \emph{LASSO} row per regression dataset for
side-by-side orientation against the external comparators;
Table~\ref{tab:matched_class} below gives the full matched-model-class
benchmark across all four datasets and three configurations and is
the authoritative source for the pipeline-then-LASSO comparison.

\begin{table}[!htbp]
\centering
\small
\caption{Pipeline-then-LASSO matched-model-class benchmark with paired $t$-tests. Same model class (LASSO for regression, $\ell_1$-Logistic for BBBP) and same 5-fold splits across all three configurations within each dataset, so the comparison isolates the effect of the upstream screening step. ``pre-LASSO feats'' is the mean number of features handed to LASSO after screening; ``final L1 feats'' is the mean number of features with non-zero LASSO coefficient. ``$t$ vs.\ full'' and ``$p$ vs.\ full'' are the paired $t$-statistic and two-sided $p$-value across the five folds (df $=4$) comparing the primary metric (RMSE for regression, ROC-AUC for BBBP) against the \texttt{full} (no-screen) row of the same dataset.}
\label{tab:matched_class}
\begin{tabular}{llcrrrr}
\toprule
Dataset & Configuration & Metric (mean $\pm$ std) & pre & final & $t$ vs.\ full & $p$ vs.\ full\\
\midrule
ESOL          & \texttt{full} (LASSO)                       & RMSE $1.536\pm 0.077$  & $30.0$         & $17.6$ & --- & ---\\
              & \texttt{pairwise} $\to$ LASSO               & RMSE $1.580\pm 0.085$  & $\phantom{0}7.4$ & $\phantom{0}7.2$ & $+3.77$ & $0.020^{*}$\\
              & \texttt{combined} $\to$ LASSO               & RMSE $1.588\pm 0.090$  & $\phantom{0}4.4$ & $\phantom{0}4.4$ & $+3.62$ & $0.022^{*}$\\
\midrule
FreeSolv      & \texttt{full} (LASSO)                       & RMSE $3.469\pm 0.400$  & $30.0$         & $12.8$ & --- & ---\\
              & \texttt{pairwise} $\to$ LASSO               & RMSE $3.465\pm 0.394$  & $\phantom{0}8.0$ & $\phantom{0}6.4$ & $-0.34$ & $0.748$\\
              & \texttt{combined} $\to$ LASSO               & RMSE $3.462\pm 0.387$  & $\phantom{0}3.2$ & $\phantom{0}3.2$ & $-0.46$ & $0.672$\\
\midrule
Lipophilicity & \texttt{full} (LASSO)                       & RMSE $1.128\pm 0.016$  & $30.0$         & $16.2$ & --- & ---\\
              & \texttt{pairwise} $\to$ LASSO               & RMSE $1.137\pm 0.013$  & $\phantom{0}8.0$ & $\phantom{0}7.0$ & $+3.45$ & $0.026^{*}$\\
              & \texttt{combined} $\to$ LASSO               & RMSE $1.220^{\,\dagger}$       & $\phantom{0}0.2$ & $\phantom{0}0.2$ & --- & ---\\
\midrule
BBBP          & \texttt{full} ($\ell_1$-Logistic)           & AUC $0.794\pm 0.019$   & $30.0$         & $11.6$ & --- & ---\\
              & \texttt{pairwise} $\to$ $\ell_1$-Logistic   & AUC $0.787\pm 0.026$   & $\phantom{0}6.0$ & $\phantom{0}6.0$ & $-1.71$ & $0.163$\\
              & \texttt{combined} $\to$ $\ell_1$-Logistic   & AUC $0.786\pm 0.028$   & $\phantom{0}4.0$ & $\phantom{0}4.0$ & $-1.56$ & $0.194$\\
\bottomrule
\end{tabular}

\vspace{0.5em}
\noindent\small{$^{*}$ significant at $\alpha=0.05$ (two-sided paired $t$-test, df $=4$). $^\dagger$ Combined screen retains zero features in $4/5$ Lipophilicity folds; the value shown is from the single fold that retained one feature. The same failure mode appears in Table~\ref{tab:ml_benchmark} and is treated as a formal limitation in \S\ref{sub:limitations}.}
\end{table}

\textbf{Headline reading (matched-model-class), with paired
$t$-tests across the five folds.}
The pipeline-then-LASSO pattern is a \emph{parsimony-versus-performance
trade-off}, not a uniformly costless reduction:

\begin{itemize}
\item On FreeSolv and BBBP, \texttt{pairwise} $\to$ LASSO / $\ell_1$-Logistic
matches \texttt{full} within CV noise: the paired $t$-test gives
$p = 0.748$ (FreeSolv) and $p = 0.163$ (BBBP), neither significant
at $\alpha = 0.05$. At these two datasets, the parsimony claim is
clean: $\sim 2\times$ fewer LASSO-active features (FreeSolv:
$6.4$ vs.\ $12.8$; BBBP: $6.0$ vs.\ $11.6$) at no detectable
performance cost.

\item On ESOL and Lipophilicity, the screen carries a
\emph{small but statistically significant} RMSE increase:
$+0.044$ RMSE on ESOL ($p = 0.020$) and $+0.009$ RMSE on
Lipophilicity ($p = 0.026$). In absolute terms these are
$\sim 2.8\%$ and $\sim 0.8\%$ degradations of the no-screen
LASSO baseline, respectively, in exchange for cutting the
LASSO-active feature count by approximately $58\%$ and $57\%$
(ESOL: $7.2$ vs.\ $17.6$; Lipo: $7.0$ vs.\ $16.2$).
\end{itemize}

The combined screen further halves the feature count on ESOL,
FreeSolv, and BBBP with a similar pattern (FreeSolv non-significant,
ESOL significantly higher RMSE), and degenerates on Lipophilicity
(same failure mode as the RF benchmark; see \S\ref{sub:limitations}).
The matched-model-class evidence is therefore: on $2$ of $4$ datasets
the parsimony claim is clean within CV noise; on the other $2$ the
trade-off is small ($<3\%$ degradation) and statistically detectable
at $5$-fold CV power. The cross-model-class qualification on
Table~\ref{tab:fs_comparison} is no longer load-bearing: the
screen-pipeline's parsimony behaviour is recoverable within model
class, with an honest characterisation of the small per-dataset
cost.

\FloatBarrier

% =====================================================================
\section{Discussion}\label{sec:discussion}
% =====================================================================

The two case studies and the cross-validated ML benchmark together
support a single methodological reading. Section~\ref{sec:wuzi_case}
shows that the pairwise diagnostic alone admits a parameter setting
(the Wuzi FreeSolv point at $(-1, -1, 1)$) that clears
$|r| < 0.95$ yet has multi-dimensional partial correlation
identically zero against the baseline on every dataset --- the
failure mode the combined diagnostic of Table~\ref{tab:decision} is
designed to catch. Section~\ref{sec:ml_benchmark} shows that the
combined screen is best read as a diagnostic rather than as a hard
filter: it flags the Wuzi failure mode but degenerates on
Lipophilicity, where the topology-to-$\log D$ partial-correlation
structure is too weak for any pairwise-pruned index to clear
$|\mathrm{pcor}| \ge 0.10$. The threshold should therefore be tied
to the QSPR utility loss the analyst is willing to tolerate, not
treated as a fixed cutoff.

Theorem~\ref{thm:bid} (Edge-Degree-Pair Basis) gives the structural
reason for the empirical redundancy ceiling observed across the
three datasets: the dimension-$10$ cap is intrinsic to the BID
construction on $\mathcal{G}_4$ rather than an artefact of any
particular dataset.

% =====================================================================
\section{Limitations and open questions}\label{sec:limitations}
% =====================================================================

We catalogue the assumptions and the open questions explicitly so
that downstream users of the pipeline can read off where its
guarantees stop and where additional methodological work is needed.

\subsection{Limitations of the present pipeline}\label{sub:limitations}

We catalogue the formal limitations as a reviewer-facing list so
the scope of every claim is bounded explicitly:

\begin{enumerate}
\item \textbf{Linear diagnostics only.} The pipeline is built on
Pearson correlation, OLS partial correlation, and PCA. Non-linear
relationships between a candidate index and the target, or between
the candidate and the baseline, are not captured. Distance
correlation, rank-based equivalents (Spearman), or
mutual-information-style estimators would close the non-linearity
gap at the cost of looser interpretability.

\item \textbf{Threshold convention.} $\tau_r = 0.95$ for the
pairwise screen and $\tau_{\mathrm{pcor}} = 0.10$ for the
partial-correlation condition are conventional but arbitrary
cutoffs and should be tuned to the QSPR-utility loss the analyst
is willing to tolerate. The qualitative conclusions of
\S\ref{sub:threshold_sensitivity} (no Wuzi grid point passes the
combined diagnostic on any dataset, at any of the nine tested
$(\tau_r, \tau_{\mathrm{pcor}})$ combinations) are stable; the
\emph{exact} pass counts are not.

\item \textbf{Alternative-family pass-counts are
single-target, not chemical-utility claims.} The $20$-candidate
experiment of \S\ref{sec:novel_case} is replicated across all four
datasets, and the qualitative conclusion (pairwise admits $15$--$17$
of $20$ candidates; combined admits $0$--$4$) is consistent across
ESOL, FreeSolv, Lipophilicity, and BBBP. However, the individual
combined passers (e.g.\ $\mathrm{InfoH\_deg}$ on ESOL and BBBP,
Triangles on FreeSolv) should still be interpreted as
``residual-signal positive on that specific target under the
$30$-index linear baseline'', not as useful QSPR descriptors in
absolute terms.

\item \textbf{Combined diagnostic should not be used as a hard
filter.} On Lipophilicity the strict combined screen
($\tau_{\mathrm{pcor}} \ge 0.10$) retains zero features in $4/5$
ML-benchmark folds; the pairwise screen alone reduces feature count
by ${\sim}4\times$ while preserving predictive performance within
CV noise, but the partial-correlation diagnostic should be read as
evidence about residual target signal rather than as an automatic
feature-selection rule. This degeneration is consistent with the
weak linear signal already present in the $30$-index baseline:
the full-baseline LASSO $5$-fold-CV $R^2$ on the three regression
datasets is $0.457$ (ESOL), $0.175$ (FreeSolv), and $0.118$
(Lipophilicity), and the per-dataset combined-pass counts on
the $20$ alternative-family candidates
(Table~\ref{tab:novel_combined}) decrease monotonically with
baseline $R^2$ across the same three datasets ($4$, $1$, $0$
combined passers respectively). Lipophilicity is therefore
not a calibration failure of the $|\mathrm{pcor}|=0.10$ threshold
in isolation; the linear topology-to-$\log D$ relationship is
sufficiently weak that once the baseline subtracts what little
linear signal exists, the residual partial correlations
naturally collapse. The Wuzi Lipophilicity result of
Table~\ref{tab:wuzi_verdict_paper1} and the Lipo ML-benchmark
column of Table~\ref{tab:ml_benchmark} both illustrate this
failure mode honestly.

\item \textbf{RandomForest is one model class; the
linear-comparator benchmark crosses model classes, but a matched-
class same-model benchmark is also reported.}
Section~\ref{sub:lasso_comparator} compares the screen-pipeline
configurations against LASSO, ElasticNet, mRMR, Boruta, and
$\ell_1$-Logistic on all four datasets. The screen-pipeline rows
use a RandomForest, and the external comparators use linear
regression (regression datasets) or logistic regression (BBBP);
that comparison is therefore cross-model-class. To eliminate this
qualification, Section~\ref{sub:pipeline_then_lasso} additionally
reports a matched-model-class same-model benchmark
(pipeline-then-LASSO vs.\ LASSO-only) on the four datasets; the
parsimony claim survives the matched comparison (within-CV-noise
RMSE / AUC at roughly half the LASSO-active features on every
dataset). A complementary benchmark with a tree-ensemble
comparator (e.g.\ XGBoost-mRMR vs.\ pipeline-then-XGBoost) is
left as future work.

\item \textbf{BID dimension bound covers $\Delta \le 4$ only.} The
Edge-Degree-Pair Basis bound (Theorem~\ref{thm:bid}) is sharp on
$\mathcal{G}_4$, which covers C / N / O / halogen and most S
chemistry; organometallics, hypervalent species, and inorganic
chemistry fall outside the dimensional cap and require a separate
analysis. Likewise, non-BID parametric families (distance-based,
spectral, mixed) are not bounded by Theorem~\ref{thm:bid} and a
structural dimension bound analogous to it for these families is
open.

\item \textbf{Single-scalar indices only.} The pipeline screens one
scalar at a time; a vector descriptor, or a small set of jointly
designed scalars, may be redundant individually but informative
collectively.
\end{enumerate}

\subsection{Open methodological questions}

The following directions are natural extensions of the framework.
We list them as open questions rather than completed work; a
formal study of any one of them would be a follow-up project.

\begin{enumerate}
\item \textbf{Extension beyond the BID family.} A structural
dimension bound analogous to Theorem~\ref{thm:bid} for
distance-based, spectral, or mixed-family parametric topological
indices is open.
\item \textbf{Threshold sensitivity.} A formal change-point
analysis of the pipeline conclusions as a function of the
thresholds $(\tau, \tau_{\mathrm{pcor}})$ would replace the
conventional cutoffs with data-driven equivalents.
\item \textbf{Non-linear diagnostics.} Distance correlation and
mutual-information-style estimators offer non-linear redundancy
measures that complement the linear Pearson / OLS / PCA core. A
joint diagnostic combining linear and non-linear redundancy
checks has not been validated against a QSPR benchmark.
\item \textbf{Web deployment.} A web-deployed version of the
pipeline that accepts a SMILES file and a candidate-index Python
function and returns the four diagnostics in seconds would lower
the activation energy for routine use.
\item \textbf{Two-stage pipeline-then-feature-selection beyond
LASSO.} Section~\ref{sub:pipeline_then_lasso} verifies the
two-stage idea for LASSO on the three regression tasks and for
$\ell_1$-logistic regression on BBBP, with paired $t$-tests across
the five CV folds: pairwise screening followed by the same
downstream model preserves performance within CV noise on
FreeSolv and BBBP, and carries a small but statistically
detectable RMSE cost ($<3\%$) on ESOL and Lipophilicity in
exchange for ${\sim}\!2\times$ fewer LASSO-active features
(Table~\ref{tab:matched_class}). What remains open is the broader
matched comparison for other downstream selectors and nonlinear /
tree-ensemble models, such as pipeline-then-XGBoost,
pipeline-then-RandomForest with Boruta- or mRMR-style selection,
or nonlinear redundancy diagnostics.
\end{enumerate}

% =====================================================================
\section{Conclusion}
% =====================================================================

We have introduced an open-source orthogonality-screening pipeline
for scalar topological indices and applied it across three
MoleculeNet QSPR regression benchmarks and one classification
benchmark. The pairwise screen preserves cross-validated
RandomForest performance within noise at roughly four-fold fewer
features, and a matched-model-class pipeline-then-LASSO benchmark
shows that comparable parsimony comes at no statistically
detectable cost on two of the four datasets and under a
three-percent RMSE cost on the other two. The Edge-Degree-Pair
Basis explains why the redundancy ceiling observed across datasets
is intrinsic to the bond-incident-degree construction rather than
an artefact of any single dataset. Requiring both a low maximum
pairwise correlation against a classical baseline and a
non-negligible partial correlation with the target operationalises
the question of whether a new index carries independent
information as a simple pre-flight check. Orthogonality, however,
is necessary but not sufficient for usefulness; non-linear
diagnostics, extensions of the basis result beyond maximum degree
four, threshold sensitivity, and a web deployment remain open.

% =====================================================================
\section*{Data and code availability}
% =====================================================================

Code, datasets, and figures are publicly available at
\url{https://github.com/Singati2/topological-index-orthogonality}.
The exact snapshot reported here is git tag
\texttt{v1.3-paper2-orthogonality}; a Zenodo archive of that tag
(with permanent DOI) will be created upon submission and the DOI
inserted in the final version of this section. The pipeline runs
in seconds per benchmark on a single laptop; adding a candidate
index requires a short Python function on the
hydrogen-suppressed molecular graph. Software dependencies and
pinned versions are listed in
\texttt{requirements.txt}.

% =====================================================================
\section*{Acknowledgements}
% =====================================================================

The authors thank colleagues for discussions of the screening
protocol and the alternative-family candidate index set. Funding
sources and institutional support to be acknowledged in the final
version of the manuscript.

% =====================================================================
\appendix
\section{Formal definitions of the 20 alternative-family candidate indices}\label{app:candidates}
% =====================================================================

For reproducibility we give the formal definition of each of the
twenty candidate indices screened in \S\ref{sec:novel_case}. Let
$G = (V, E)$ be a simple connected hydrogen-suppressed molecular
graph with $n = |V|$ vertices, $m = |E|$ edges, degree sequence
$\{d_v\}_{v\in V}$, geodesic distance matrix $D = (d_{uv})$,
adjacency matrix $A$, eccentricity $\varepsilon(v) = \max_u d_{uv}$,
and adjacency-spectrum / Laplacian-spectrum eigenvalues
$\{\lambda_i\}$ / $\{\mu_i\}$ respectively. All twenty definitions
are implemented in
\texttt{src/novel\_candidates.py} of the companion repository.

\paragraph{Information-theoretic (5).}
With $H(\{c_i\}) = -\sum_i (c_i / \sum_j c_j)\log(c_i / \sum_j c_j)$
the Shannon entropy of a count distribution:
\begin{align*}
\mathrm{InfoH\_deg}(G) &= H\bigl(\{|\{v : d_v = k\}|\}_k\bigr), \\
\mathrm{InfoH\_dist}(G) &= H\bigl(\{|\{(u,v) : d_{uv} = k,\, u<v\}|\}_k\bigr), \\
\mathrm{InfoH\_spec}(G) &= -\sum_{i} p_i \log p_i,\ p_i = \lambda_i^2 / \textstyle\sum_j \lambda_j^2, \\
\mathrm{InfoH\_ecc}(G) &= H\bigl(\{|\{v : \varepsilon(v) = k\}|\}_k\bigr), \\
\mathrm{Bonchev\_Id}(G) &= W\log_2 W - \sum_{k \ge 1} k\,d_k \log_2(k\,d_k),
\end{align*}
where in the last line $d_k = |\{(u,v) : d_{uv}=k,\,u<v\}|$ and
$W = \sum_k k\,d_k$ is the Wiener index.

\paragraph{Centrality-based (5).}
With standard NetworkX-style normalised centralities:
\begin{align*}
\mathrm{Btw\_sum}(G) &= \sum_{v \in V} c_B(v),\quad c_B(v) = \!\!\sum_{s \ne t \ne v}\!\!
\frac{\sigma_{st}(v)}{\sigma_{st}}, \\
\mathrm{Cls\_sum}(G) &= \sum_{v} (n-1)\Bigl(\sum_{u \ne v} d_{uv}\Bigr)^{-1}, \\
\mathrm{Eig\_sum}(G) &= \sum_{v} x_v,\ x \text{ the (positive) Perron eigenvector of } A, \\
\mathrm{Harm\_cent}(G) &= \sum_{v}\sum_{u \ne v}\frac{1}{d_{uv}}, \\
\mathrm{Btw\_var}(G) &= \mathrm{Var}_v\bigl(c_B(v)\bigr).
\end{align*}

\paragraph{Spectral-hybrid (3).}
With $\mathrm{diam}(G) = \max_{u,v} d_{uv}$ and Estrada index
$\mathrm{EE}(G) = \sum_i e^{\lambda_i}$:
\begin{align*}
\mathrm{AlgConn}(G) &= \mu_2(G), \text{ second-smallest Laplacian eigenvalue}, \\
\mathrm{Rho}\!\times\!\mathrm{Dia}(G) &= \rho(G)\cdot\mathrm{diam}(G),\ \rho(G) = \max_i|\lambda_i|, \\
\mathrm{EE}\!\times\!\mathrm{W}(G) &= \log\bigl(1+\mathrm{EE}(G)\bigr)\cdot\log\bigl(1+W(G)\bigr).
\end{align*}

\paragraph{Motif / structural (4).}
\begin{align*}
\mathrm{Triangles}(G) &= \tfrac{1}{6}\operatorname{tr}(A^3), \\
\mathrm{MeanClust}(G) &= \frac{1}{n}\sum_{v} \frac{2 t(v)}{d_v(d_v-1)},\ \,t(v) = \text{triangles at } v, \\
\mathrm{FourCyc}(G) &= \tfrac{1}{8}\bigl(\operatorname{tr}(A^4) - 2\!\sum_v d_v(d_v - 1) - 2m\bigr), \\
\mathrm{CutVerts}(G) &= |\{v \in V : G - v \text{ is disconnected}\}|.
\end{align*}

\paragraph{Eccentricity-based (3).}
\begin{align*}
\mathrm{TotEcc}(G) &= \sum_{v \in V} \varepsilon(v), \\
\mathrm{Radius}(G) &= \min_{v \in V} \varepsilon(v), \\
\mathrm{MeanEcc}(G) &= \frac{1}{n}\sum_{v \in V} \varepsilon(v).
\end{align*}

None of these twenty candidates is claimed as a newly invented
index. The information-theoretic family follows Bonchev and
Trinajsti\'c~\cite{Gutman2013Degree, TodeschiniConsonniHandbook};
the centrality family follows the standard graph-theory
literature; the motif and eccentricity families are textbook
graph invariants; the spectral-hybrid combinations are
deliberately constructed novel pairings to test whether
non-standard product indices escape the redundancy ceiling. They
appear here purely as screening targets to test the generality
of the combined diagnostic of \S\ref{sec:protocol}, and the
verdict reported in \S\ref{sec:novel_case} is a statement about
this particular candidate set under the tested baselines and
thresholds, not a claim that any individual candidate is or is
not a useful QSPR descriptor.

\bibliographystyle{wileyNJD-Harvard}
\bibliography{paper2_references}

\end{document}
